\section{Stationary-shell frame and spectral compression}
\label{sec:shell-frame}

The shell transform is built before the stationary restriction. This order is
important: the full transform is exactly Parseval, and only the explicitly
omitted nonstationary modes create a defect.

\subsection{Critical packet lattice, translation reserve, and Fourier trim}
Put
\[
R_{\rm pkt}:=\frac{C_{\rm pkt}}2\log L,
\qquad
L_{\rm pkt}=L+2R_{\rm pkt}.
\]
Use the full flat packet Fourier interval
\[
0\le u\le L_{\rm pkt},
\]
whose central physical zeta band is
$R_{\rm pkt}\le u\le R_{\rm pkt}+L$. By Proposition~\ref{prop:trivial-outer-twist}, the invariant contour source has
no nontrivial outer-denominator translation. The logarithmic reserve is used
only for the packet-dimension slack and the smooth transition caps; it does not
encode an arithmetic truncation.

For exact lattice centres
\[
\tau_k=\tau_0+\frac{2\pi k}{L_{\rm pkt}}
\]
put
\begin{equation}
S_v(u):=L_{\rm pkt}^{-1/2}\sum_kv_ke^{i\tau_ku}.
\label{eq:lattice-synthesis-shell}
\end{equation}
Fourier orthogonality gives the exact isometry
\begin{equation}
\boxed{
\int_0^{L_{\rm pkt}}|S_v(u)|^2du=\|v\|_2^2.
}
\label{eq:lattice-synthesis-isometry}
\end{equation}
Recall $W_T,X_T$ from \eqref{eq:stable-log-core}. Let $C_T$ be the stable-core
restriction of Definition~\ref{def:PT}. On packet coefficients its positive
Gram is
\[
\langle P_Wv,v\rangle
:=\int_{R_{\rm pkt}}^{R_{\rm pkt}+W_T}|S_v(u)|^2du
=\|C_T S_v\|_2^2.
\]
This is the same core restriction used in $P_T$, $A_T^G$, and $\Gfr_T$; it is
not an additional packet-coordinate deletion.

\begin{lemma}[Stable-core packet frame]
\label{lem:stable-core-frame}
For every admissible packet-coordinate compression there is a subspace
$V_{\rm core}$ such that
\begin{equation}
\boxed{
\codim V_{\rm core}\ll T\log L=o(N_T),
\qquad
\|C_TS_v\|_2^2\ge c_{\rm core}\|v\|_2^2
\quad(v\in V_{\rm core}),
}
\label{eq:stable-core-frame}
\end{equation}
with a fixed $c_{\rm core}>0$. The corresponding bound for the diagonal
multiplier $v_\nu\mapsto\tau_\nu^{-1}v_\nu$ holds with the right side multiplied
by $T^{-2}$.
\end{lemma}

\begin{proof}
For exact lattice centres, \eqref{eq:lattice-synthesis-isometry} is Parseval on
an interval of length $L_{\rm pkt}$. Its complement to the stable core has
length
\[
L_{\rm pkt}-W_T=O(\log L).
\]
If $D_{\rm tail}$ is the Gram of restriction to that complement, then each
coordinate vector contributes $O(\log L/L)$ to its trace. Hence on a packet
fibre of dimension $O(TL)$,
\[
\tr D_{\rm tail}\ll T\log L.
\]
Removing the spectral subspace of $D_{\rm tail}$ for eigenvalues at least
$1/4$ costs $O(T\log L)$ dimensions and leaves the core Gram at least $3/4$ of
the full flat Gram. Denote this good packet subspace by $V_0$.

By Proposition~\ref{prop:trivial-outer-twist} there is only the identity outer
translation. For the primitive estimate intersect $V_0$ with
$D_\tau^{-1}V_0$, where
$D_\tau v=(\tau_\nu^{-1}v_\nu)_\nu$. Since $D_\tau$ is invertible and
$T\le\tau_\nu\le2T+O(1)$, this preimage has the same codimension as $V_0$, while
\[
\|D_\tau v\|_2^2\asymp T^{-2}\|v\|_2^2.
\]
Thus the intersection has codimension $O(T\log L)$ and gives both lower bounds
in \eqref{eq:stable-core-frame}. Moving every centre by at most $T^{-100}$
changes these finite Grams by $o(1)$ in operator norm, by the estimate in the
proof of Lemma~\ref{lem:packet-riesz}; the fixed lower bound therefore persists.
\end{proof}

\subsection{Logarithmic unitary and exact local Parseval}
On the retained physical band write
\[
\widetilde u:=u-R_{\rm pkt},
\qquad x=e^{\widetilde u},
\]
and define
\begin{equation}
F_v(x):=x^{-1/2}S_v(R_{\rm pkt}+\log x),
\qquad 1\le x\le X_T.
\label{eq:log-unitary}
\end{equation}
Since $dx=x\,d\widetilde u=x\,du$,
\begin{equation}
\boxed{
\int_1^{X_T}|F_v(x)|^2dx
=\int_{R_{\rm pkt}}^{R_{\rm pkt}+W_T}|S_v(u)|^2du
=\langle P_Wv,v\rangle.
}
\label{eq:log-unitary-norm}
\end{equation}
We use an explicit fixed-overlap partition by length-one windows. Choose
$\vartheta\in C_c^\infty((-1/2,1/2))$ which is positive on
$[-1/4,1/4]$, put $c_j=j/2$, and define
\[
D(x):=\left(\sum_{j\in\mathbb Z}\vartheta(x-c_j)^2\right)^{1/2},
\qquad
\psi_j(x):=\frac{\vartheta(x-c_j)}{D(x)}.
\]
Because every real number lies within distance $1/4$ of some $c_j$, the
$1/2$-periodic function $D$ is smooth and bounded above and below by positive
constants. With
\[
I_j=[J_j,J_j+1],
\qquad J_j:=c_j-\frac12,
\]
we therefore have
\begin{equation}
\boxed{
\supp\psi_j\subset I_j,
\qquad
\sum_{j\in\mathbb Z}\psi_j(x)^2=1,
\qquad
\|\psi_j^{(r)}\|_\infty\ll_r1.
}
\label{eq:smooth-shell-partition}
\end{equation}
Extend $F_v$ by zero outside $[1,X_T]$ and retain shells with
$\psi_j\not\equiv0$ on $(1,X_T)$. The $O(1)$ shells meeting the two core
endpoints remain in the Parseval decomposition; their endpoint jets are treated
by Proposition~\ref{prop:boundary-endpoint-rank}.

\begin{lemma}[Global triangular metric and shell-local primitives]
\label{lem:global-local-triangular}
Let $g$ be Hilbert-valued on $[0,W_T]$ with $\int_0^{W_T}g=0$, and put
\[
R(u):=\int_0^u g(v)\,dv,
\qquad f(x):=x^{-1/2}g(\log x).
\]
Then $R(0)=R(W_T)=0$. Extend $f$ and
$\mathcal R(x):=R(\log x)$ by zero outside $[1,X_T]$. After the normalization
$r=u/L$, the flat triangular primitive form is
\begin{equation}
\mathcal Q_{\rm glob}(g)
:=\frac{2}{L^2}\int_0^{W_T}\|R(u)\|^2\,du
=\frac{2}{L^2}\int_1^{X_T}
\left\|\int_1^x y^{-1/2}f(y)\,dy\right\|^2\frac{dx}{x}.
\label{eq:global-triangular-log}
\end{equation}
Let $\mathcal J_T$ denote these retained shells. For $j\in\mathcal J_T$ define
\[
q_j(x):=\psi_j(x)x^{-1/2}f(x),
\qquad
Q_j(x):=\int_{J_j}^xq_j(y)\,dy,
\]
and
\begin{equation}
\mathcal Q_{\rm loc}(f)
:=\frac{2}{L^2}\sum_{j\in\mathcal J_T}
\int_{I_j}\|Q_j(x)\|^2\frac{dx}{x}.
\label{eq:local-triangular-shell}
\end{equation}
Then there is an absolute constant $C_\triangle$ such that
\begin{equation}
\boxed{\mathcal Q_{\rm loc}(f)\le
C_\triangle\mathcal Q_{\rm glob}(g).}
\label{eq:global-local-triangular-comparison}
\end{equation}
The estimate is Hilbert-valued. If, in addition,
$\int_{I_j}q_j=0$ for every shell and output component, then $Q_j$ vanishes at both
endpoints and may be used in the two-variable primitive identity of
Lemma~\ref{lem:two-variable-primitive}.
\end{lemma}

\begin{proof}
The zero extension of $\mathcal R$ is absolutely continuous because
$R(0)=R(W_T)=0$, and $\mathcal R'=x^{-1/2}f$ almost everywhere. Thus
\[
Q_j(x)=\psi_j(x)\mathcal R(x)
-\int_{J_j}^x\psi_j'(y)\mathcal R(y)\,dy,
\]
because $\psi_j$ vanishes at the shell endpoints. The windows have uniformly
bounded first derivative and bounded overlap. On the physical part of a
length-one shell, $x^{-1}\asymp(1+J_j)^{-1}$, so Cauchy--Schwarz gives
\[
\int_{I_j}\|Q_j(x)\|^2\frac{dx}{x}
\ll\int_{I_j}\|\mathcal R(x)\|^2\frac{dx}{x}.
\]
Summing in $j$ proves \eqref{eq:global-local-triangular-comparison} and is
unchanged for Hilbert-valued functions. The final assertion follows directly
from the definition of $Q_j$.
\end{proof}

\begin{corollary}[Tensor stability of the global/local Gram comparison]
\label{cor:tensor-global-local-triangular}
Let $\mathcal H$ be any Hilbert space, with no restriction on its finite
compression dimension.  If $g:[0,W_T]\to\mathcal H$ and the corresponding
$f,Q_j$ satisfy the hypotheses of
Lemma~\ref{lem:global-local-triangular}, then
\[
\mathcal Q_{\rm loc}(f)\le
C_\triangle\mathcal Q_{\rm glob}(g)
\]
with the same constant $C_\triangle$, independently of $\dim\mathcal H$.
Consequently the comparison remains valid after adjoining any packet-independent
finite source, refinement, Cauchy, or carrier output fibre by Hilbert tensor
product or orthogonal direct sum.
\end{corollary}

\begin{proof}
The proof of Lemma~\ref{lem:global-local-triangular} uses only the Hilbert norm,
Bochner integration, Cauchy--Schwarz, and the bounded overlap of the scalar
windows $\psi_j$.  Each step is dimension free.  Applying the same argument in
$\mathcal H$, or equivalently first on finite orthogonal sums and then by
completion, gives the assertion.
\end{proof}

Define
\begin{equation}
Q_{j,n}(v):=\int_{I_j}\psi_j(x)F_v(x)e^{-2\pi inx}\,dx.
\label{eq:shell-coeff}
\end{equation}
Because each shell has length one, the integer exponentials are an orthonormal
basis and Parseval gives
\[
\sum_{n\in\mathbb Z}|Q_{j,n}(v)|^2
=\|\psi_jF_v\|_2^2.
\]
Summing in $j$,
\begin{equation}
\boxed{
\|\mathsf Q_{\rm full}v\|_2^2
=\langle P_Wv,v\rangle.
}
\label{eq:full-shell-isometry}
\end{equation}
No stationary approximation has entered this identity.

\subsection{Stationary band and nonstationary leakage}
For one packet centre $\tau$ and one shell Fourier mode $n$, the phase is
\[
\Phi_{\tau,n}(x)=\tau\log x-2\pi nx,
\qquad
\Phi'_{\tau,n}(x)=\frac\tau x-2\pi n.
\]
For $T\le\tau\le2T$ and $x\in I_j$, a stationary point can occur only for
\begin{equation}
\mathcal I_j
:=\left[\frac{T}{2\pi(J_j+1)},
    \frac{2T}{2\pi J_j}\right].
\label{eq:stationary-interval}
\end{equation}
Choose the shell margin
\begin{equation}
R_j:=L^B\left(1+\frac{\sqrt T}{J_j}\right)
\label{eq:shell-margin}
\end{equation}
and let $\mathcal I_j^+$ contain the integers at distance at most $R_j$ from
$\mathcal I_j$. Let $\mathsf Q_{\rm shell}$ retain only $(j,n)$ with
$n\in\mathcal I_j^+$ and let $\mathsf Q_{\rm out}$ be its complement.

Fix a packet basis vector $e_k$. On a nonboundary shell put
$D=\dist(n,\mathcal I_j)\ge R_j$. Then
\[
|\Phi'|\ge2\pi D,
\quad |\Phi^{(r)}|\ll_rT/J_j^r,
\quad |A_{j,k}^{(r)}|\ll_rL_{\rm pkt}^{-1/2}J_j^{-1/2}L^{Cr}J_j^{-r},
\quad \frac{T/J_j^2}{D^2}\ll L^{-2B}.
\]
Repeated integration by parts gives
\begin{equation}
\boxed{|Q_{j,n}(e_k)|
\ll_NL_{\rm pkt}^{-1/2}J_j^{-1/2}L^{C_N}(1+D)^{-N}.}
\label{eq:nonstat-shell}
\end{equation}
Thus the squared nonboundary contribution is
\[
\ll_N\frac{L^{2C_N}}{L_{\rm pkt}}
\sum_jJ_j^{-1}R_j^{1-2N}.
\]
For a boundary shell, split instead at the physical endpoint
$b\in\{1,X_T\}$. Repeating the same integration by parts $N$ times gives a
finite endpoint sum plus an interior remainder. The endpoint sum is a linear
combination of
\begin{equation}
\boxed{
\partial_x^r\!\bigl(\psi_j(x)F_v(x)\bigr)\big|_{x=b},
\qquad 0\le r<N,
}
\label{eq:boundary-shell-jets}
\end{equation}
with scalar coefficients depending on the Fourier mode and the phase; the
interior remainder satisfies the same arbitrary-order bound as
\eqref{eq:nonstat-shell}. Thus the physical endpoint contributes only finitely
many packet jet functionals, while every nonstationary interior term remains in
the regular integration-by-parts class.

The shell enlargement itself has only $o(N_T)$ extra rows. Indeed
\begin{align*}
\sum_jR_j
&\ll X_TL^B+\sqrt T\,L^B\log X_T\\
&\ll TL^{B-A_0}+\sqrt T\,L^{B+1}=o(N_T)
\end{align*}
when $A_0>B+3$.

\begin{proposition}[Terminal logarithmic tail]
\label{prop:terminal-log-tail}
Let $X_{\rm full}:=e^L=T/(2\pi)$. On the complementary logarithmic band
\[
W_T<u\le L,
\qquad
X_T<x\le X_{\rm full},
\]
the vertical right-edge source and its adjoint decompose as
\begin{equation}
\boxed{
A_T^{\rm tail}
=E_T^{\rm tail}+E_T^{\rm bdry,tail}+R_T^{\rm tail},
}
\label{eq:terminal-log-tail}
\end{equation}
where $E_T^{\rm bdry,tail}$ is the tail-side subfamily of the split endpoint
jets at $x=X_T$ and is included, exactly once, in the global boundary matrix
$E_T^{\rm bdry}$ of Proposition~\ref{prop:boundary-endpoint-rank}. Moreover
\begin{equation}
\boxed{
\rank E_T^{\rm tail}
\ll T\log L=o(N_T),
}
\label{eq:terminal-log-tail-rank}
\end{equation}
and, on $V_{\rm core}\cap V_{\rm mean}$,
\begin{equation}
\boxed{
\|\Gfr_T^{-1/2}R_T^{\rm tail}\Gfr_T^{-1/2}\|_{S_2}^2=o(N_T).
}
\label{eq:terminal-log-tail-bounds}
\end{equation}
The same estimates include the core--tail cross terms.
\end{proposition}

\begin{proof}
Use the same unit-shell Fourier transform on $X_T<x\le X_{\rm full}$, but put
only the integers at distance at most $1$ from $\mathcal I_j$ into the tail
stationary range. Its total number of rows is
\[
\sum_{X_T\ll J_j\ll X_{\rm full}}
\left(1+\frac{T}{J_j}\right)
\ll T+T\log\frac{X_{\rm full}}{X_T}
\ll T\log L.
\]
By Proposition~\ref{prop:trivial-outer-twist} there is no arithmetic fibre
multiplicity. Every stationary tail term and every core--tail cross term
therefore factors through a space of dimension $O(T\log L)$, which proves the
rank bound independently of its direct-source coefficients.  For the full
exact--model residual, Lemmas~\ref{lem:direct-carrier-square} and
\ref{lem:stationary-symbol} give a uniform polylogarithmic coefficient-square
and symbol-derivative loss; this is absorbed into the logarithmic constants
below.

For the complementary modes let $D=\dist(n,\mathcal I_j)>1$. Since
$J_j\ge X_T=T/(2\pi L^{A_0})$,
\[
\frac{T}{J_j^2}\ll\frac{L^{2A_0}}{T}.
\]
The integration-by-parts argument of \eqref{eq:nonstat-shell}, now with the
fixed tail margin $1$, gives for every fixed $N$
\[
|Q_{j,n}(e_k)|
\ll_N L_{\rm pkt}^{-1/2}J_j^{-1/2}
\left(\frac{L^{C_N}}{J_j}\right)^N(1+D)^{-N}.
\]
Squaring this estimate and summing first in $n$ gives a bounded
$\sum_{D>1}(1+D)^{-2N}$. Since the retained unit shells satisfy
$J_j\ge X_T\asymp T L^{-A_0}$,
\[
\sum_{J_j\ge X_T}J_j^{-1-2N}\ll X_T^{-2N}.
\]
After summing over the $O(TL)$ packet coordinates, with no additional
arithmetic multiplicity by Proposition~\ref{prop:trivial-outer-twist}, we
obtain the sufficient absolute estimate
\begin{equation}
\boxed{
\|R_T^{\rm tail}\|_{S_2}^2
\ll T^{1-2N}L^{C_{N,A_0}}.
}
\label{eq:terminal-tail-absolute-HS}
\end{equation}
The same estimate holds for the nonstationary core--tail cross terms, since the
core analysis map has fixed operator norm up to a power of $L$.
Lemma~\ref{lem:core-reference-floor} gives on
$V_{\rm core}\cap V_{\rm mean}$
\[
\|\Gfr_T^{-1/2}\|_{\op}^2\ll T^2L^{C_G}.
\]
Consequently
\[
\|\Gfr_T^{-1/2}R_T^{\rm tail}\Gfr_T^{-1/2}\|_{S_2}^2
\ll T^{5-2N}L^{C_{N,A_0}+2C_G}=o(N_T),
\]
because the fixed order $N$ chosen in Section~\ref{sec:retained} is in
particular larger than $3$. Endpoint jets at the physical endpoint
$x=X_{\rm full}$ factor through finitely many evaluation functionals and are
included in $E_T^{\rm tail}$. The tail-side split jets at $x=X_T$ are precisely
$E_T^{\rm bdry,tail}$ and belong only to the global boundary matrix
$E_T^{\rm bdry}$ of Proposition~\ref{prop:boundary-endpoint-rank}. This proves
\eqref{eq:terminal-log-tail} and \eqref{eq:terminal-log-tail-bounds}.
\end{proof}

\section{Macroscopic \texorpdfstring{one-$\chi$}{one-chi} shifts have low edge rank}
\label{sec:edge-rank}

For a direct stationary one-$\chi$ source on a unit shell write
\[
\sum_mW_{j,m}a_{j,m}(y;p)e(-my),\qquad
a_{j,m}=\sum_r\widehat a_{j,m}(r;p)e(ry).
\]
Here $m\in\mathcal I_j$. The sideband $r$ has effective carrier $q=m-r$, since
$e(ry)e(-my)=e(-qy)$. If $P_j$ projects onto $\mathcal I_j^+$ and
$(S_qb)_n=b_{n+q}$, its compression is $P_jS_qP_j$. Only
$|r|\le R_j^{\rm sb}$ enters the edge matrix; the complement is regular
class~(vii).

Write
\[
A_j:=\frac{T}{2\pi(J_j+1)},
\qquad
B_j:=\frac{2T}{2\pi J_j}.
\]
If $m\in\mathcal I_j$, then $P_jS_mP_j$ can have nonzero range only on the
common low-edge interval
\[
A_j-R_j\le n\le B_j+R_j-A_j.
\]
Thus all forward shifts share one edge space $E_j^-$ with
\begin{equation}
\boxed{
\dim E_j^-
\ll B_j-2A_j+2R_j+1
\ll \frac{T}{J_j^2}+R_j+1.
}
\label{eq:edge-shell-dimension}
\end{equation}
The adjoints lie in a high-edge space $E_j^+$ of the same dimension. Therefore,
for arbitrary scalar weights $W_{j,m}$,
\begin{equation}
\rank\left(
\sum_{m\in\mathcal I_j}W_{j,m}P_jS_mP_j+\mathrm{adjoint}
\right)
\ll
\frac{T}{J_j^2}+R_j+1.
\label{eq:edge-shell-rank}
\end{equation}
The same conclusion holds with coefficients acting on an auxiliary common fibre
only when they preserve this shell-coordinate edge range; no such extra fibre is
present in the invariant source by Proposition~\ref{prop:trivial-outer-twist}.
For the parameter $\epsilon_L=L^{-\sigma}$ fixed in
Section~\ref{sec:retained}, the low Fourier sidebands of a smooth stationary
symbol replace $m$ by the effective shift $m-r$ with
$|r|\le\epsilon_Ld_j$, where $d_j\asymp T/J_j$. Enlarging $E_j^\pm$ by this amount therefore changes the
right side of \eqref{eq:edge-shell-rank} only by $O(\epsilon_Ld_j)$. Now
\[
\sum_j\frac{T}{J_j^2}\ll T,
\qquad
\sum_jR_j
\ll T L^{B-A_0}+\sqrt T\,L^{B+1}=o(T),
\]
and
\[
\sum_j\epsilon_Ld_j
\ll \epsilon_LT\log X_T
\ll T L^{1-\sigma}.
\]
Thus the first two geometric edge contributions are $O(T)$. The low-sideband
enlargement contributes $T L^{1-\sigma}=o(T)$ once $\sigma>2$. Hence
\begin{proposition}[Uniform rank bound for low-sideband stationary shifts]
\label{prop:edge-rank}
Let $E_T^{\rm shift}$ be any finite sum of the low-sideband stationary
compressed shifts described above, with base carrier indices in the intervals
$\mathcal I_j$ and sidebands $|r|\le R_j^{\rm sb}$. Then
\begin{equation}
\boxed{
\rank E_T^{\rm shift}
\ll T=o(N_T).
}
\label{eq:edge-rank-detailed}
\end{equation}
The estimate depends only on the geometric range of the retained low-sideband
shifts and is independent of their coefficients or signs. It also holds for
norm-convergent sums of such shifts, since every partial sum has range in the
same fixed finite-dimensional edge space.  In the invariant-source final route
this proposition is applied to the genuinely direct stationary carrier
families, in particular the entire exact--model residual
\eqref{eq:direct-exact-model-residual} with its full norm-summable HLP
hierarchy. The complementary sidebands are not part of this rank assertion.
\end{proposition}

\begin{proposition}[Stable-core endpoint terms have finite rank]
\label{prop:boundary-endpoint-rank}
Let $E_T^{\rm bdry}$ be the sum of the endpoint terms produced by the fixed
$N$ integrations by parts on the $O(1)$ shell windows meeting
$b\in\{1,X_T\}$, including both one-sided contributions at the split
$x=X_T$. Denote the contribution from the tail side of this split by
$E_T^{\rm bdry,tail}$, as in \eqref{eq:terminal-log-tail}. Then
\begin{equation}
\boxed{
\rank E_T^{\rm bdry}
\ll L^{C_{\rm bdry}}=o(N_T)
}
\label{eq:boundary-endpoint-rank}
\end{equation}
for a fixed exponent $C_{\rm bdry}$.
\end{proposition}

\begin{proof}
For each boundary shell, fixed local-source label $\alpha$, endpoint $b$, and
$0\le r<N$, let
\[
(J_Tv)_{b,r,\alpha}
:=\partial_x^r\!\bigl(\psi_j(x)F_{v,\alpha}(x)\bigr)\big|_{x=b}.
\]
Equation~\eqref{eq:boundary-shell-jets} shows that every endpoint contribution
with one boundary leg factors as $J_T^*BQ$ for some output operator $B$ and
some companion analysis $Q$; the reflected term factors as $Q^*B^*J_T$.
Hence each Hermitian pair has rank at most $2\rank J_T$. There are two physical
endpoints, $O(1)$ boundary shells, and the derivative order $N$ is fixed. The fixed local-source labels contribute at most $L^{C_{\rm bdry}}$
coordinates. Thus
\[
\rank J_T\ll L^{C_{\rm bdry}}=o(N_T).
\]
\end{proof}

Define $E_T^{\rm edge}$ to be the sum of the low-sideband stationary shifts of
the entire direct exact--model residual \eqref{eq:direct-exact-model-residual}
and of the other genuinely direct stationary sources, together with the
boundary-shell matrix $E_T^{\rm bdry}$ and the terminal logarithmic tail matrix
$E_T^{\rm tail}$.  Proposition~\ref{prop:edge-rank} applies to the full
norm-convergent HLP sum because its range bound is independent of coefficients
and signs.  The auxiliary outer-excess term is zero by
Proposition~\ref{prop:trivial-outer-twist}; the horizontal packet-centre collars
are handled only by $V_{\rm end}$.  Propositions~\ref{prop:edge-rank},
\ref{prop:boundary-endpoint-rank}, and \ref{prop:terminal-log-tail} give
\begin{equation}
\boxed{\rank E_T^{\rm edge}\ll T\log L=o(N_T).}
\label{eq:total-edge-rank}
\end{equation}
